\subsection{Skeleton-trace optimization}
\label{subsec:app-skeleton-trace-optimization}

\paragraph{Exact interface.}
The target is the exact trace functional $L_S$ from
Definition~\ref{def:trace-length}; the compact body consequence is
Theorem~\ref{thm:common-skeleton-count}.

\paragraph{Variables and feasible set.}
The C-triangle variables range over the signed CE1/CE2 domain of
\zcref{prop:signed-center-normal-form}.  A V triangle ranges over
the closed unit equilateral triangles containing its distinguished vertex,
with the displayed reach and radial-incidence restrictions in the lemmas
below.

\paragraph{Objective.}
For each indicated class, maximize $L_S(T)$.

\paragraph{Exact result.}
On the four feasible classes below, the exact sufficient caps are,
respectively,
\[
 L_S(T_C)<\frac32,\qquad
 L_S(T_i)<\frac32,\qquad
 L_S(T_i)\le2,\qquad
 L_S(T_i)<\frac32.
\]

\paragraph{Solution.}
The following four lemmas give the complete calculations for those caps.

\begin{lemma}[Center skeleton cap]
\label{lem:center-skeleton-cap}
If $T_C$ is a CE1 or CE2 C triangle and contains exactly one radial
midpoint, then
\[
 L_S(T_C)<\frac32.
\]
\end{lemma}

\begin{proof}
Normalize the unique midpoint as $M_0$ and the positive boundary trace to
$e_{0,1}$.  Use the signed variables of
\zcref{prop:signed-center-normal-form}; in particular
\[
 W=1-R,\qquad E=\sqrt{1-RW},\qquad
 \eta=1-E,\qquad P=E(1-E).
\]
The positive trace condition is
\begin{equation}
 \alpha+W\delta<P.
 \label{eq:center-skeleton-domain}
\end{equation}
The two possible boundary contributions are
\[
 \ell_{01}=W+\delta-\frac{\eta+\alpha+\delta}{R},
 \qquad
 \ell_{50}=\frac{[P-(R\alpha+\delta)]_+}{W}.
\]
The six radial contributions are bounded by
\[
\begin{array}{c|cccccc}
 i&0&1&2&3&4&5\\ \hline
 \mathcal H^1(T_C\cap r_i)&
 E-\alpha-\delta&\delta/R&\delta&
 \min\{\alpha/R,\delta/W\}&\alpha&\alpha/W.
\end{array}
\]
Adding gives
\[
 L_S(T_C)\le K_R+\mathcal E(\alpha,\delta),
\]
where
\[
 K_R=W+E-\frac\eta R
\]
and
\[
 \mathcal E(\alpha,\delta)
 =-\frac\alpha R+\frac\alpha W+\delta
 +\min\left\{\frac\alpha R,\frac\delta W\right\}
 +\frac{[P-(R\alpha+\delta)]_+}{W}.
\]
Splitting only according to the minimum and the positive part gives
\begin{equation}
 \mathcal E(\alpha,\delta)\le\frac PW.
 \label{eq:center-skeleton-E-bound}
\end{equation}
For example, if $\alpha/R\le\delta/W$, the expression without the positive
part is $\alpha/W+\delta$; when the positive part is active, its excess over
$P/W$ is $\alpha-R\delta/W\le0$.  The reflected order gives excess
$\delta-W\alpha/R\le0$.  If the positive part vanishes,
\eqref{eq:center-skeleton-domain} gives the same bound strictly.

Put $A=1+R-R^2$.  Then
\[
 \frac32-K_R-\frac PW
 =\frac{1+A-2EA}{2RW}.
\]
Both sides of $1+A>2EA$ are positive, and
\[
 (1+A)^2-4A^2E^2
 =R^2W^2(5+4R-4R^2)>0.
\]
Thus $K_R+P/W<3/2$.
\end{proof}

\begin{lemma}[Positive-support skeleton cap]
\label{lem:positive-support-skeleton-cap}
Let $T_i$ be the closure of an original open V triangle, so
$V_i\in\mathrm{int}(T_i)$.  If
\[
 \Haus(T_i\cap r_j)>0
 \quad\text{for some }j\in\{i-1,i+1\},
\]
then
\[
 L_S(T_i)<\frac32.
\]
\end{lemma}

\begin{proof}
Normalize to $i=0$ and translate the hexagon by $V_0$:
\[
 \widetilde H=V_0+H,\qquad \widetilde O=V_0,\qquad W_j=V_0+V_j.
\]
Then $W_3=O$, $W_2=V_1$, and $W_4=V_5$.  The five original skeleton
components on which a unit triangle containing $V_0$ can have
positive-length trace are
\[
 [V_0,V_1],\ [V_0,V_5],\ [O,V_1],\ [O,V_0],\ [O,V_5].
\]
All five belong to the skeleton $\widetilde S$ of $\widetilde H$.  Since
$V_0=\widetilde O$ is interior to $T_0$, the triangle is a legitimate center
triangle for $\widetilde H$.  Its positive trace on $r_1$ or $r_5$ is a positive
boundary-edge trace of $\widetilde H$, so it is CE1 or CE2 by
Proposition~\ref{prop:ce-classification}.  Proposition
\ref{prop:unique-center-midpoint} and Lemma~\ref{lem:center-skeleton-cap}
therefore give
\[
 L_{\widetilde S}(T_0)<\frac32.
\]
A nonlocal component of $S$ is more than unit distance from the interior
point $V_0$, apart from zero-measure endpoints.  Hence
$L_S(T_0)\le L_{\widetilde S}(T_0)$.
\end{proof}

\begin{lemma}[No-support skeleton cap]
\label{lem:no-support-skeleton-cap}
If a V triangle has $A_i+B_i\le1$ and
\[
 \Haus(T_i\cap r_{i-1})=\Haus(T_i\cap r_{i+1})=0,
\]
then
\[
 L_S(T_i)\le2.
\]
\end{lemma}

\begin{proof}
Its boundary contribution is $A_i+B_i\le1$ by
\eqref{eq:vertex-boundary-locality}.  For $P=sV_j\in r_j\setminus\{O\}$,
\[
 \lVert V_i-P\rVert^2
 =1+s^2-2s\cos\frac{(j-i)\pi}{3}>1
\]
when $j\notin\{i-1,i,i+1\}$.  The traces on $r_{i-1}$ and $r_{i+1}$ have
zero length by hypothesis, so only $r_i$ can contribute positive length.
This adds at most one.
\end{proof}

For the remaining cap we use one specialized part of the exact local
equilateral enclosure criterion.  In local directions $u,v$ meeting at
$120^\circ$, suppose a closed unit triangle contains
\[
 0,\quad au,\quad bv,\quad c(u+v),
\]
with $0\le a\le b\le1$ and $a+b>1$.  The support-line calculation gives the
selected supercritical cell
\begin{equation}
 \begin{aligned}
 &a^2+ab+b^2\le1,\qquad c\le\frac12,\\
 &(a^2-1)c^2+(2ab^2+b)c+(b^4-b^2)\le0,
 \end{aligned}
 \label{eq:supercritical-admissible-cell}
\end{equation}
where $c$ lies on the connected component containing $c=0$.  To see the
selector directly, the quadratic is negative at $c=0$, positive at
$c=1/2$, and opens downward.  For the middle assertion, four times its value
at $1/2$ is increasing in $a$ and, at $a=1-b$, equals
$b^2(2b-1)^2$; it is therefore positive when $a+b>1$.
Hence $c\le1/2$ selects its smaller-root component.  Formula
\eqref{eq:supercritical-admissible-cell} follows by
fixing in the support sum the three active contacts $\{B\}$,
$\{B,c(u+v)\}$, and $\{A\}$, and squaring the positive unit-side equation.

\begin{lemma}[Supercritical skeleton cap]
\label{lem:supercritical-skeleton-cap}
If $A_i+B_i>1$, then
\[
 L_S(T_i)<\frac32.
\]
\end{lemma}

\begin{proof}
By Proposition~\ref{prop:vertex-classification}, a V triangle is Vd0,
Vd1, Vd2, or T3-like.  Proposition~\ref{prop:vd-corner-normal-form} makes
every Vd1/Vd2 V triangle nonsupercritical, and
Lemma~\ref{lem:t3-nonsupercritical} does the same for T3-like V triangles.  Hence a
supercritical triangle is Vd0 and, by definition,
$\Haus(T_i\cap r_{i-1})=\Haus(T_i\cap r_{i+1})=0$.  Thus, after normalizing
and writing $a=A_i$, $b=B_i$, the skeleton trace has length $a+b+c$, where
$c$ is the reach on $r_i$.
Reflect so that $a\le b$, and put
\[
 s=a+b,\qquad \delta=b-a,\qquad c_0=\frac32-s.
\]
The diameter constraint gives
\[
 1<s\le\frac2{\sqrt3},\qquad
 0\le\delta\le\sqrt{4-3s^2}.
\]
Let $P(c)$ denote the quadratic in
\eqref{eq:supercritical-admissible-cell}.  Substitution and multiplication by
$16$ give
\[
 \begin{aligned}
 16P(c_0)=Q(s,\delta)
 &=\delta^4+8\delta^3s-6\delta^3+14\delta^2s^2-18\delta^2s+5\delta^2\\
 &\quad-8\delta s^3+30\delta s^2-34\delta s+12\delta\\
 &\quad+s^4-6s^3-19s^2+60s-36.
 \end{aligned}
\]
Its derivative factors as
\[
 \frac{\partial Q}{\partial\delta}
 =2(2\delta+4s-3)
 \bigl(\delta^2+\delta(4s-3)+(s-1)(2-s)\bigr)\ge0.
\]
Therefore
\[
 Q(s,\delta)\ge Q(s,0)
 =(s-1)(s^3-5s^2-24s+36)>0.
\]
Indeed, the cubic is decreasing on $[1,2/\sqrt3]$ and at its right endpoint
equals $88/3-136\sqrt3/9>0$.  Hence $P(c_0)>0$.  Since
\eqref{eq:supercritical-admissible-cell} selects the smaller-root component,
$P(c)\le0$ forces $c<c_0$.  Thus $a+b+c<3/2$.
\end{proof}

\paragraph{Body consequence.}
These four optima supply the four skeleton rows used in
Theorem~\ref{thm:common-skeleton-count}.

\subsection{Vd1/Vd2 corner optimization}
\label{subsec:app-vd-corner-optimization}

\paragraph{Exact interface.}
The public outputs are the Vd1/Vd2 boundary row of
Theorem~\ref{thm:boundary-trace-table} and the branch in
Proposition~\ref{prop:ce2-vd2-midpoint-length}.

\paragraph{Variables and feasible set.}
Let $T$ be the closure of an original Vd1 or Vd2 triangle at $V_0$.
Use the temporary coordinates
\[
 X=V_0+x(V_5-V_0)+y(V_1-V_0).
\]
Thus $V_0=(0,0)$, $V_5=(1,0)$, $V_1=(0,1)$, $O=(1,1)$, and
$\lVert(x,y)\rVert^2=x^2+y^2-xy$.

\paragraph{Objective.}
Determine the exact half-plane family and maximize the boundary sum, first
without and then with an adjacent-midpoint containment constraint.

\paragraph{Exact result.}
The exact half-plane and ray descriptions are
\eqref{eq:vd-normal-form} and \eqref{eq:vd-ray-intervals}.  Their optimized
consequences are
\[
 A+B<\frac12,
 \qquad
 M_{i-1}\in T_i\ \text{or}\ M_{i+1}\in T_i
 \Longrightarrow A_i+B_i<\frac13.
\]

\paragraph{Solution.}

\begin{proposition}[Vd1/Vd2 corner normal form]
\label{prop:vd-corner-normal-form}
Let $T$ be the closure of an original Vd1 or Vd2 triangle at $V_0$, and let
$A,B$ be its exact reaches toward $V_5,V_1$.  There is a unique $t>0$ such
that, with $d=\sqrt{t^2+t+1}$,
\[
 T=\left\{(x,y):
 \begin{aligned}
 x-(t+1)y&\le A,\\
 ty-(t+1)x&\le tB,\\
 tx+y&\le d-A-tB
 \end{aligned}\right\}.
 \label{eq:vd-normal-form}
\]
The raw traces, in coordinates measured from the indicated outer vertex, are
\[
 \begin{array}{c|c}
 r_0&0\le s\le\dfrac{d-A-tB}{t+1}\\[6pt]
 r_1&\dfrac{t(1-B)}{t+1}\le s\le\dfrac{d-A-tB-1}{t}\\[8pt]
 r_5&\dfrac{1-A}{t+1}\le s\le d-A-tB-t.
 \end{array}
 \label{eq:vd-ray-intervals}
\]
Intersecting with $[0,1]$ gives the actual traces.  Moreover,
\begin{equation}
 A+B<\frac12.
 \label{eq:vd-half-cap}
\end{equation}
\end{proposition}

\begin{proof}
A Vd1/Vd2 triangle has one vertex $W$ outside $H$ and two vertices $U,Z$
in $H$.  Since $V_0$ is an interior convex combination of them, both
coordinates of $W$ are negative.  The two axis endpoints $(A,0),(0,B)$ lie
on distinct sides incident with $W$.  Put $p=U-W$, $q=Z-W$.  Both vectors
have positive coordinates and satisfy
\[
 \lVert p\rVert=\lVert q\rVert=\lVert p-q\rVert=1.
\]
The order in which the two sides cross the positive axes selects the
$60^\circ$ rotation
\[
 q=(p_x-p_y,p_x),\qquad p_x>p_y>0.
\]
Writing $t=(p_x-p_y)/p_y$ gives uniquely
\[
 p=\frac{(t+1,1)}d,\qquad q=\frac{(t,t+1)}d.
\]
The two side lines through the axis endpoints and the parallel third side
are exactly those in \eqref{eq:vd-normal-form}.  Substitution of
$(s,s),(s,1),(1,s)$ gives \eqref{eq:vd-ray-intervals}.

If the $r_1$ trace has positive length, its lower endpoint is smaller than
its upper endpoint.  Clearing denominators gives
\[
 (t+1)A+tB<(t+1)(d-1)-t^2.
\]
Since the left side is at least $t(A+B)$,
\[
 A+B<\frac{(t+1)(d-1)-t^2}{t}<\frac12.
\]
For the last inequality, both sides are positive and
\[
 \left(t^2+\frac{3t}{2}+1\right)^2
 -(t+1)^2(t^2+t+1)=\frac{t^2}{4}>0.
\]
Reflection proves the same conclusion when the supported arm is $r_5$.
\end{proof}

\begin{lemma}[Vd2 adjacent-midpoint cap]
\label{lem:vd2-neighbor-midpoint-cap}
If an open role $U_i$ has Vd2 closure $T_i=\overline{U_i}$ and $T_i$
contains $M_{i-1}$ or $M_{i+1}$, then its
exact boundary reaches satisfy
\[
 A_i+B_i<\frac13.
\]
\end{lemma}

\begin{proof}
Normalize to $i=0$ and write $A=A_i$, $B=B_i$.  In
Proposition~\ref{prop:vd-corner-normal-form}, reflection sends
$(t,A,B,M_1,M_5)$ to $(1/t,B,A,M_5,M_1)$, so assume $t\ge1$ and put
$U=A+tB$.

If $M_5=(1,1/2)$ is contained, the third half-plane in
\eqref{eq:vd-normal-form} gives
\[
 A+B\le U\le d-t-\frac12\le\sqrt3-\frac32<\frac13;
\]
the middle expression decreases for $t\ge1$.

If $M_1=(1/2,1)$ is contained, the second and third half-planes give
\[
 B\ge\frac{t-1}{2t},\qquad
 A+B\le S(t):=d-t-\frac1{2t}.
\]
Positive-length contact with $r_5$ gives, by
\eqref{eq:vd-ray-intervals},
\[
 A+B<R(t):=\frac{2(t+1)d-3t^2-t-2}{2t}.
\]
Direct differentiation shows that $S$ is increasing and $R$ decreasing on
$[1,\infty)$.  For example, the sign calculation for $S'$ reduces to
\[
 (2t+1)^2t^4-(t^2+t+1)(2t^2-1)^2
 =(t+1)(t^3+3t^2-1)>0,
\]
while $R'<0$ follows from $d>t+1/2$.  Splitting at $t=4/3$ yields
\[
 \begin{aligned}
 1\le t\le\frac43&:\quad
 A+B\le S(4/3)=\frac{8\sqrt{37}-41}{24}<\frac13,\\
 t\ge\frac43&:\quad
 A+B<R(4/3)=\frac{7\sqrt{37}-39}{12}<\frac13.
 \end{aligned}
\]
The final inequalities follow respectively from
$8\sqrt{37}<49$ and $7\sqrt{37}<43$.
\end{proof}

\paragraph{Body consequence.}
The first result supplies the Vd1/Vd2 row of
Theorem~\ref{thm:boundary-trace-table}; the second supplies
Proposition~\ref{prop:ce2-vd2-midpoint-length}.

\subsection{Boundary-trace optimization}
\label{subsec:app-boundary-trace-optimization}

\paragraph{Exact interface.}
The objective is the body functional $L_{\partial H}$, and the public result is
Theorem~\ref{thm:boundary-trace-table}.

\paragraph{Variables and feasible set.}
The center ranges over the three exact classes determined by $N_E(T_C)$.
Each V triangle ranges over its exact normalized V type and its actual maximal
reaches $(A_i,B_i)$.

\paragraph{Objective.}
Maximize $L_{\partial H}$ separately on every row of the public table.

\paragraph{Exact result.}
The optimum in each feasible class is the corresponding entry of
\eqref{eq:boundary-trace-table}.

\paragraph{Solution.}
The signed center normal form gives both positive center traces from one
formula, and the remaining rows follow from the exact V-triangle constraints.

\begin{lemma}[Boundary-trace table calculation]
\label{lem:app-boundary-trace-table-calculation}
For closures of original open triangles, the following bounds hold:
\[
 \begin{array}{c|c|c}
 \text{triangle}&\text{hypothesis}&L_{\partial H}\\ \hline
 T_C&\mathrm{CE0}&0\\
 T_C&\mathrm{CE1}&\le\dfrac{\sqrt3}{2}-\dfrac34\\[3pt]
 T_C&\mathrm{CE2}&<\dfrac12\\[3pt]
 T_i&\mathrm{Vd0},\ A_i+B_i\le1&\le1\\
 T_i&\mathrm{Vd0},\ A_i+B_i>1&\le\dfrac2{\sqrt3}\\[3pt]
 T_i&\mathrm{Vd1/Vd2}&<\dfrac12\\
 T_i&\text{T3-like}&<1.
 \end{array}
 \label{eq:app-boundary-trace-table}
\]
\end{lemma}

\begin{proof}
For every V triangle,
\begin{equation}
 L_{\partial H}(T_i)=A_i+B_i.
 \label{eq:vertex-boundary-locality}
\end{equation}
Indeed, the two incident traces are initial intervals of these lengths.  A
relative-interior point of any nonincident edge is more than unit distance
from $V_i$, whereas $V_i$ is interior to the diameter-one triangle $T_i$.
This excludes any other positive-length boundary trace.

The CE0 C-triangle row is the definition.  The other two C-triangle rows are
the class specializations of the common signed formula in
\zcref{cor:signed-center-boundary}.

For a nonsupercritical Vd0 V triangle, \eqref{eq:vertex-boundary-locality} is enough.
For a supercritical V triangle, its incident endpoints are at squared distance
$A_i^2+A_iB_i+B_i^2\le1$, and hence
\[
 \frac34(A_i+B_i)^2\le1.
\]
This proves the $2/\sqrt3$ cap.  The Vd1/Vd2 V triangle is
\eqref{eq:vd-half-cap}.  Finally, the Type II calculation in
Proposition~\ref{prop:t3-translation} gives for an original T3-like triangle
\[
 A_i+B_i=z-\alpha<z<1
\]
because $0<t<1$ and $\alpha>0$.
\end{proof}

\paragraph{Body consequence.}
The seven exact maxima are the rows of
Theorem~\ref{thm:boundary-trace-table}.
